Put \(p=n-1\). For rational profiles \(u_i=-v(I_i)/q_i\), allowing unit ideals, define
\[
B_{\boldsymbol u}(v(J))=\frac{e_R(I_1,\ldots,I_p,J)}{q_1\cdots q_p}.
\]
For \(p=0\) this means \(e_R(J)\) with empty denominator one. N1 proves input-function independence, symmetry, addition, positive rational scaling, and ideal-test order.

Every rational ideal-difference test is \((v(J)-v(K))/d\). Equality with \((v(J')-v(K'))/d'\) gives
\[
v(J^{d'}(K')^d)=v(K^{d'}(J')^d)
\]
on ALL supports. Proper products have equal integral closures; equality with a unit forces both products to be units. The power/product laws prove equal proposed values, rational linearity, and all test-relation independence. A nonnegative difference gives closure inclusion and hence functional nonnegativity by n1; unit cases are immediate or impossible. Thus the whole rational lattice has a positive functional.

The constant \(1=v(\mathfrak a)\) belongs to that lattice. Positivity on \(r\pm f\), rational \(r>\|f\|_\infty\), gives \(|B(f)|\le\|f\|_\infty B(1)\) after decreasing \(r\). G331 supplies uniform density in real \(C(X)\). Norm-continuous extension is real-linear by rational scalar approximation. It is positive: for \(h\ge0\), approximate by lattice \(f_j\), add rational \(r_j>\|h-f_j\|\) tending to zero, and pass \(B(f_j)\ge-r_jB(1)\). Compact-Hausdorff Riesz representation gives a unique finite positive Radon measure \(\mu_{\boldsymbol u}\), of mass \(e_R(I_1,\ldots,I_p,\mathfrak a)/\prod q_i\). No metrizability is assumed.

Set \(E_J=e_R(\mathfrak a^{[p]},J)\), \(E_{\mathfrak a}=e_R(\mathfrak a)\). Comparing profiles in \([-M,0]\) to \(-M=-Mv(\mathfrak a)\), first rationally and then by rational upper limits, gives
\[
0\le B_{\boldsymbol u}(v(J))\le M^pE_J,\qquad
\mu_{\boldsymbol u}(X)\le M^p E_{\mathfrak a}.
\]
If two rational tuples have coordinate distances at most \(\epsilon\), use \(u_i-\epsilon\le w_i\) and the reverse inequalities. Expanding after adding the legitimate negative constant in every slot gives
\[
|B_{\boldsymbol u}(v(J))-B_{\boldsymbol w}(v(J))|
\le\Delta_p(M,\epsilon)E_J,\qquad
\Delta_p(M,\epsilon)=(M+\epsilon)^p-M^p. \tag{U}
\]
Each nonempty subset \(S\) of shifted slots contributes a nonnegative term at most \(\epsilon^{|S|}M^{p-|S|}E_J\); their sum is the stated polynomial. Rational upper errors remove restrictions on bounds. Set \(\Delta_0=0\), \(M^0=1\).

Approximate \(C\) tuples uniformly by rational primary profiles, using G331. A common finite bound and (U) make every ideal moment Cauchy and independent of approximation. Differences define the lattice limit. Positivity and mass bounds persist; Riesz gives the measure. Symmetry, addition and real nonnegative scaling pass to limits. If \(u\le w\), approximants within \(\eta\) satisfy \(u_{\rm app}-2\eta\le w_{\rm app}\); (U) removes this shift, proving order. The same bounds and (U) therefore hold throughout \(C\). Convergence of ideal moments and common mass bounds imply weak convergence for every continuous test by uniform lattice approximation. This proves uniform weak continuity for nets as well.

For bounded \(u_i\in D\), take ANY actual decreasing \(C\)-representations \(g_{i,\alpha_i}\downarrow u_i\). If \(u_i\ge-M\), all terms are already in \([-M,0]\). On the finite product index set, ideal moments increase and are bounded. Their suprema, and differences of these limits, define a positive norm-bounded lattice functional. Riesz gives its measure; density and mass bounds prove weak convergence of representing measures.

To prove independence, let \(f_{i,\beta_i}\downarrow u_i\) be another representation. Fix \(\boldsymbol\beta\) and \(\epsilon>0\). The directed open sets \(\{g_{i,\alpha_i}<f_{i,\beta_i}+\epsilon\}\) cover compact \(X\). Finite subcovers and directedness select one \(\alpha_i\) for each coordinate valid everywhere. Thus
\[
B_{\boldsymbol f_\beta}(v(J))
\le B_{\boldsymbol g_\alpha-\epsilon}(v(J))
\le\sup_{\boldsymbol\alpha}B_{\boldsymbol g_\alpha}(v(J))+\Delta_p(M,\epsilon)E_J.
\]
Decrease \(\epsilon\) to zero, take the supremum over \(\boldsymbol\beta\), and reverse the roles. Moments and measures coincide. The same cover argument for a fixed \(C\)-majorant \(h_i\ge u_i\) proves
\[
B_{\boldsymbol u}(v(J))
=\sup\{B_{\boldsymbol h}(v(J)):h_i\in C,\ h_i\ge u_i\}. \tag{M}
\]
Every such majorant belongs to \([-M,0]\); representing terms prove the reverse bound. This yields ideal-moment order on bounded \(D\): the majorants of a larger tuple form a smaller set. It asserts no order of measures.

Addition, scaling and symmetry on \(D\) follow by actual finite product \(C\)-representations and scalar moment limits. Sums of representing nets decrease to the required sums. Thus the negative shift expansion and (U) hold on bounded \(D\) as well. Uniform convergence INSIDE this established domain gives weak convergence. Increasing closure is not used.

For actual decreasing-net continuity let \(u_{i,\lambda}\in D\) decrease with common bound \([-M,0]\) to a SPECIFIED \(u_i\in D\). Order bounds the increasing moments above by the target. Fix a \(C\)-majorant tuple \(h\) and \(\epsilon>0\). The \(D\) functions are upper semicontinuous, so \(\{u_{i,\lambda}<h_i+\epsilon\}\) are directed open covers. Compactness selects a common \(\lambda\), giving
\[
B_{\boldsymbol h}(v(J))\le B_{\boldsymbol u_\lambda-\epsilon}(v(J))
\le B_{\boldsymbol u_\lambda}(v(J))+\Delta_p(M,\epsilon)E_J.
\]
Take the supremum over \(\lambda\), decrease \(\epsilon\) to zero, and take the supremum in (M). Every target ideal moment is recovered, and density and mass bounds prove weak convergence on every continuous test. Limit membership in \(D\) is a hypothesis. No nested-net flattening, arbitrary-infimum closure, or unbounded-energy limit is asserted.
